\documentclass[runningheads]{llncs}

% Định dạng Springer CCIS/LNCS cho SOICT 2026, đồng bộ với bản tiếng Anh.
% Không thêm geometry, fancyhdr, titlesec, setspace hoặc gói font tùy chỉnh:
% lớp llncs kiểm soát khổ giấy, lề, tiêu đề và kiểu chữ theo yêu cầu.
\usepackage[T5]{fontenc}
\usepackage[utf8]{inputenc}
\usepackage[vietnamese]{babel}
\usepackage{amsmath,amssymb}
\usepackage{graphicx}
\usepackage{booktabs}
\usepackage{url}
\usepackage[hidelinks]{hyperref}

\renewcommand{\abstractname}{Tóm tắt}
\providecommand{\keywordname}{}
\renewcommand{\keywordname}{\textbf{Từ khóa:}}

\begin{document}

\title{LegalSCM-RAG: Truy xuất đường dẫn nhân quả và xác minh cấu trúc chọn lọc cho hỏi đáp pháp luật đa bước tiếng Việt}
\titlerunning{LegalSCM-RAG}

% SOICT 2026 sử dụng phản biện single-blind, vì vậy giữ tên tác giả và cơ quan.
\author{Nguyễn Thị Nhung, Nguyễn Quốc Tuấn, Nguyễn Thị Kim Anh}
% \authorrunning{N. T. Nhung và cộng sự}

\institute{Đại học Bách khoa Hà Nội, Hà Nội, Việt Nam\\
\email{ }}

\maketitle
% SOICT 2026 yêu cầu bản nộp PDF không có số trang.
\pagestyle{empty}

\begin{abstract}
Các mô hình ngôn ngữ lớn và kiến trúc sinh tăng cường truy xuất vẫn gặp hạn chế
với câu hỏi pháp luật yêu cầu suy luận điều kiện--hệ quả nhiều bước. Bài báo đề
xuất \textbf{LegalSCM-RAG}, kết hợp truy xuất đường dẫn nhân quả, suy luận quy
tắc quy phạm ba giá trị và xác minh cấu trúc chọn lọc. Các sự kiện pháp lý được
biểu diễn bằng hai kênh hỗ trợ đúng và sai, trong khi các ma trận incidence thưa
cho thân và đầu quy tắc mã hóa phép hội trong một quy tắc và hỗ trợ tuyển thay
thế giữa nhiều quy tắc. Chúng tôi cũng phân biệt tính đạt được sau khi xóa nút,
chỉ được dùng như một chẩn đoán tô-pô, với can thiệp cấu trúc $do(X_M=0)$.
Phép can thiệp này cố định sự kiện trung gian $M$ ở trạng thái sai, thay thế cơ
chế tương ứng và suy luận lại điểm bất động của hệ quy tắc.

Trên 250 câu hỏi được sinh từ đồ thị Bộ luật Hình sự Việt Nam, Full LegalSCM đạt
độ chính xác quyết định 96,80\%, so với 88,80\% của Causal Path RAG và 55,20\%
của Vanilla RAG. Mức tăng theo cặp 8,0 điểm so với Causal Path RAG có ý nghĩa
thống kê (CI 95\%: 4,4--12,0; $p$ hiệu chỉnh Holm $=0,0011$), trong khi Answer
F1, ROUGE-L và Citation F1 không khác biệt có ý nghĩa. Các đối chứng closed-book
đạt 4,40\% khi cho phép abstention và 9,20\% khi buộc quyết định nhị phân; các
kết quả này phản ánh hành vi dưới output contract kiểm chứng thay vì năng lực
pháp luật tổng quát. Hệ thống thực thi 7 phép $do(X_M=0)$, nhưng Causal Path RAG
đã trả lời đúng cả 7 câu hỏi tương ứng. Vì vậy, kết quả hỗ trợ hiệu quả của toàn
bộ cơ chế xác minh chọn lọc, còn lợi ích riêng của can thiệp cấu trúc cần được
xác lập thêm.

\keywords{Sinh tăng cường truy xuất \and Hỏi đáp đa bước
\and Logic ba giá trị \and Lập trình logic đại số tuyến tính
\and Can thiệp cấu trúc \and Hỏi đáp pháp luật tiếng Việt}
\end{abstract}



\section{Giới thiệu}
\label{sec:introduction}

Các mô hình ngôn ngữ lớn (Large Language Models -- LLMs) đạt kết quả tốt trên
nhiều bài toán hỏi đáp. Tuy nhiên, một hệ thống hỏi đáp pháp luật đáng tin cậy
cần cả căn cứ luật thành văn phù hợp và chuỗi suy luận có thể truy vết, liên kết
các điều kiện, sự kiện trung gian và hệ quả pháp lý
\cite{chalkidis2020legalbert,guha2023legalbench}. Sinh tăng cường truy xuất
(Retrieval-Augmented Generation -- RAG) cung cấp bằng chứng bên ngoài trước
bước sinh \cite{lewis2020retrieval,gao2024ragsurvey}; tuy nhiên, truy xuất ngữ
nghĩa truyền thống thường xem các đoạn văn là độc lập và không biểu diễn tường
minh các quan hệ điều kiện--hệ quả nhiều bước.

Chúng tôi đề xuất \textbf{LegalSCM-RAG}, kết hợp truy xuất đường dẫn nhân quả
với xác minh cấu trúc phụ thuộc câu hỏi. Đồ thị biểu diễn các quan hệ phụ thuộc
quy phạm được trích xuất từ luật thành văn, không phải hiệu ứng nhân quả được
ước lượng từ dữ liệu quan sát. Tầng quy tắc sử dụng ngữ nghĩa điểm bất động ba
giá trị \cite{kleene1952metamathematics,fitting1985kripke} và biểu diễn đại số
tuyến tính thưa dựa trên các nghiên cứu về logic programming trong không gian
vector và tensor \cite{sato2020threevalued,sakama2021tensor}. Các ma trận
incidence cho thân và đầu quy tắc mã hóa phép hội giữa các literal trong cùng
một quy tắc và hỗ trợ tuyển thay thế giữa các quy tắc có cùng hệ quả.

Bộ xác minh chuyển claim trực tiếp tới phép kiểm tra cạnh với hai đầu mút đã
được ground, câu hỏi thông thường tới suy luận quy tắc factual, và câu hỏi loại
bỏ sự kiện trung gian tường minh tới can thiệp cứng $do(X_m=0)$. Chúng tôi phân
biệt can thiệp này với tính đạt được sau khi xóa nút: kiểm tra tính đạt được
trong $G^{-m}$ chỉ là một chẩn đoán tô-pô có giới hạn, trong khi
$do(X_m=0)$ cố định sự kiện trung gian $m$ ở trạng thái sai, thay thế cơ chế
cấu trúc tương ứng và suy luận lại điểm bất động. Selective commit gate chỉ
chấp nhận quyết định của bộ xác minh khi các điều kiện về grounding, độ phủ
đường, confidence và tính nhất quán trạng thái được thỏa mãn.

Các đóng góp chính của nghiên cứu gồm:
\begin{itemize}
    \item một framework truy xuất và xếp hạng các đường suy luận nhiều bước
    trên quan hệ điều kiện--hệ quả trong luật thành văn;
    \item một formulation quy tắc ba giá trị dựa trên đại số tuyến tính, hỗ trợ
    phép hội, các cơ chế tuyển thay thế và can thiệp cấu trúc tường minh; và
    \item một đánh giá có kiểm soát trên năm cấu hình, sử dụng kiểm định ý nghĩa
    theo cặp cùng các audit về route, commit, abstention và intervention.
\end{itemize}

Trên 250 câu hỏi được sinh từ đồ thị Bộ luật Hình sự Việt Nam, Full LegalSCM
đạt độ chính xác quyết định 96,80\%, so với 88,80\% của Causal Path RAG và
55,20\% của Vanilla RAG. Mức tăng theo cặp 8,0 điểm so với Causal Path RAG có
ý nghĩa thống kê (CI 95\%: $4{,}4$--$12{,}0$; $p$ hiệu chỉnh Holm
$=0{,}0011$), trong khi Answer F1, ROUGE-L và Citation F1 không khác biệt có ý
nghĩa. Độ chính xác của closed-book vẫn thấp khi cho phép abstention (4,40\%)
và khi buộc quyết định nhị phân (9,20\%). Các kết quả này phản ánh hành vi dưới
output contract định hướng xác minh khi không có bằng chứng pháp luật được truy
xuất, thay vì năng lực pháp luật tổng quát của LLM. Mặc dù hệ thống thực thi
bảy can thiệp $do(X_m=0)$, Causal Path RAG cũng trả lời đúng cả bảy câu hỏi
tương ứng. Vì vậy, kết quả hỗ trợ hiệu quả tổng thể của cơ chế xác minh chọn
lọc, còn lợi ích riêng của can thiệp cấu trúc cần được đánh giá thêm.

% =========================================================
% II. RELATED WORK
% =========================================================
\section{Nghiên cứu liên quan}
\label{sec:related_work}

\subsection{Sinh tăng cường truy xuất và hỏi đáp đa bước}

Sinh tăng cường truy xuất (Retrieval-Augmented Generation -- RAG) kết hợp mô
hình ngôn ngữ với bằng chứng bên ngoài được truy xuất trước bước sinh câu trả
lời \cite{lewis2020retrieval,gao2024ragsurvey}. Các bộ truy xuất dense như DPR
biểu diễn câu hỏi và đoạn văn trong một không gian embedding chung
\cite{karpukhin2020dense}. Các nghiên cứu tiếp theo cải thiện việc tổng hợp
bằng chứng, truy xuất thích ứng và tổ chức tri thức theo cấu trúc phân cấp
\cite{izacard2021fid,asai2024selfrag,sarthi2024raptor}. Tuy nhiên, phần lớn hệ
thống RAG vẫn cung cấp bằng chứng chủ yếu dưới dạng các đoạn văn độc lập thay vì
một cấu trúc suy luận tường minh.

Các benchmark hỏi đáp đa bước như HotpotQA \cite{yang2018hotpotqa} và MuSiQue
\cite{trivedi2022musique} cho thấy việc truy xuất các đoạn văn liên quan riêng
lẻ là chưa đủ khi câu trả lời phụ thuộc vào nhiều bước suy luận liên kết. Hạn
chế này đặc biệt quan trọng trong hỏi đáp pháp luật, nơi một kết luận có thể phụ
thuộc vào chuỗi điều kiện, sự kiện trung gian, ngoại lệ và hệ quả pháp lý.

\subsection{RAG dựa trên đồ thị và quan hệ nhân quả}

Truy xuất dựa trên đồ thị tổ chức bằng chứng thông qua các quan hệ giữa thực
thể, sự kiện hoặc đoạn văn. GraphRAG khai thác cấu trúc đồ thị để tổng hợp thông
tin từ các nguồn liên quan \cite{edge2024graphrag}, trong khi Think-on-Graph
duyệt đồ thị một cách tường minh để xây dựng các đường suy luận có thể truy vết
\cite{sun2024thinkongraph}. CausalRAG đưa quan hệ nhân quả vào quá trình truy
xuất và suy luận \cite{wang2025causalrag}, còn CC-RAG xây dựng các chuỗi nhân
quả để hỗ trợ suy luận nhiều bước có cấu trúc \cite{parekh2025ccrag}.

Bài toán của chúng tôi khác ở chỗ các cạnh của đồ thị biểu diễn quan hệ phụ
thuộc điều kiện--hệ quả mang tính quy phạm được trích xuất từ văn bản luật.
Chúng được sử dụng để tổ chức suy luận pháp lý và không nên được hiểu là hiệu
ứng nhân quả thống kê được ước lượng từ dữ liệu quan sát. Ngoài ra, tính đạt
được trên đồ thị không tự nó xác định ngữ nghĩa của các quy tắc pháp luật liên
quan.

\subsection{Hỏi đáp pháp luật và suy luận phản thực}

Các benchmark và mô hình chuyên biệt cho pháp luật, bao gồm LexGLUE, CaseHOLD,
LegalBench và LEGAL-BERT, cho thấy cả đặc trưng riêng của ngôn ngữ pháp luật và
độ khó của suy luận pháp lý có cấu trúc
\cite{chalkidis2022lexglue,zheng2021casehold,guha2023legalbench,
chalkidis2020legalbert}. Các nghiên cứu này tạo động lực cho những hệ thống kết
hợp truy xuất căn cứ pháp luật với cấu trúc suy luận tường minh và có thể audit.

Nghiên cứu về hỏi đáp đa bước phản thực cũng cho thấy mô hình có thể sinh câu
trả lời đúng nhưng dựa trên một chuỗi suy luận rời rạc hoặc không phù hợp
\cite{guo2023counterfactual}. Tuy nhiên, trong mô hình nhân quả cấu trúc, một
phép can thiệp không tương đương với việc xóa nút khỏi đồ thị
\cite{pearl2009causality}. Phép xóa nút chỉ kiểm tra tính đạt được về mặt tô-pô,
trong khi $do(X_m=0)$ cố định một biến và thay thế cơ chế cấu trúc của biến đó
trước khi các biến hậu duệ được suy luận lại. Sự phân biệt này cần thiết khi
đánh giá vai trò của một sự kiện pháp lý trung gian.

\subsection{Lập trình logic ba giá trị và đại số tuyến tính}

Logic ba giá trị phân biệt trạng thái đúng, sai và chưa xác định hoặc chưa biết
\cite{kleene1952metamathematics}. Ngữ nghĩa Kripke--Kleene của Fitting cung cấp
cách diễn giải điểm bất động cho các chương trình logic trong điều kiện thông
tin không đầy đủ \cite{fitting1985kripke}. Sato và cộng sự mã hóa ngữ nghĩa
chương trình logic hai và ba giá trị trong không gian vector
\cite{sato2020threevalued}, trong khi Sakama và cộng sự sử dụng biểu diễn ma
trận và tensor cho các chương trình logic xác định, tuyển và chuẩn
\cite{sakama2021tensor}.

Các nghiên cứu này tạo cơ sở cho formulation incidence thưa đối với thân và đầu
quy tắc của chúng tôi. Phép hội được biểu diễn bằng nhiều literal trong cùng
một quy tắc, còn hỗ trợ tuyển thay thế được biểu diễn bằng các quy tắc kích hoạt
độc lập có cùng hệ quả. Formulation cung cấp một cách nhìn đại số tuyến tính cho
hệ quy tắc, nhưng runtime đánh giá các quy tắc thưa tương ứng bằng phép suy luận
ký hiệu tới điểm bất động thay vì thực hiện phép nhân ma trận hoặc tensor dày.

\subsection{Khoảng trống nghiên cứu}

Các phương pháp RAG hiện có tập trung vào truy xuất ngữ nghĩa, các phương pháp
dựa trên đồ thị tập trung vào kết nối cấu trúc, còn lập trình logic đại số tuyến
tính cung cấp ngữ nghĩa quy tắc hình thức. Tuy nhiên, vẫn còn ít nghiên cứu kết
hợp các hướng này trong một framework hỏi đáp pháp luật đồng thời:

\begin{itemize}
    \item truy xuất và xếp hạng các đường suy luận pháp luật nhiều bước;
    \item biểu diễn trạng thái pháp lý chưa đầy đủ bằng ngữ nghĩa ba giá trị;
    \item xử lý phép hội trong một quy tắc và các cơ chế tuyển thay thế giữa
    nhiều quy tắc;
    \item phân biệt phép xóa nút tô-pô với can thiệp cấu trúc; và
    \item commit có chọn lọc các quyết định đã được xác minh trước bước sinh.
\end{itemize}

LegalSCM-RAG giải quyết khoảng trống này bằng cách tích hợp truy xuất đường dẫn
nhân quả, suy luận điểm bất động ba giá trị, xác minh cấu trúc phụ thuộc câu hỏi
và commit gate có khả năng audit cho bài toán hỏi đáp pháp luật đa bước tiếng
Việt.

% =========================================================
% III. PROBLEM FORMULATION
% =========================================================
\section{Phát biểu bài toán}
\label{sec:problem_formulation}

Chúng tôi xét bài toán hỏi đáp pháp luật nhiều bước, trong đó một kết luận có
thể yêu cầu liên kết các điều kiện luật định, sự kiện trung gian và hệ quả pháp
lý. Tri thức pháp luật được biểu diễn bằng đồ thị có hướng $G=(V,E)$ và tập quy
tắc $\mathcal{R}$. Mỗi nút biểu diễn một sự kiện pháp lý đã được chuẩn hóa, còn
mỗi cạnh tóm tắt một hoặc nhiều quy tắc điều kiện--hệ quả. Đồ thị hỗ trợ truy
xuất và phân tích tô-pô, trong khi $\mathcal{R}$ xác định ngữ nghĩa suy luận.

Với câu hỏi $q$, hệ thống xây dựng các đường suy luận ứng viên
\begin{equation}
\mathcal{P}(q)=\{P_1,\ldots,P_K\},\qquad
P_i=v_0\xrightarrow{r_1}\cdots\xrightarrow{r_h}v_h,
\end{equation}
và chọn đường chính
\begin{equation}
P^*=\arg\max_{P_i\in\mathcal{P}(q)}S(P_i,q),
\end{equation}
trong đó $S$ kết hợp độ phù hợp ngữ nghĩa và cấu trúc.

\subsection{Ngữ nghĩa ba giá trị dựa trên đại số tuyến tính}

Mỗi sự kiện được gắn với một biến
\begin{equation}
X_i\in\{1,0,\bot\},
\end{equation}
lần lượt biểu diễn \textsc{true}, \textsc{false} và \textsc{unknown}. Trạng
thái unknown thể hiện hỗ trợ bị thiếu hoặc xung đột, không phải một giá trị số
nằm giữa 0 và 1. Formulation này dựa trên ngữ nghĩa điểm bất động ba giá trị
\cite{kleene1952metamathematics,fitting1985kripke} và các mở rộng trong không
gian vector \cite{sato2020threevalued}.

Trạng thái sự kiện được mã hóa bằng hai kênh hỗ trợ:
\begin{equation}
(t_i,f_i)=
\begin{cases}
(1,0), & X_i=1,\\
(0,1), & X_i=0,\\
(0,0), & X_i=\bot.
\end{cases}
\label{eq:two_channel_state}
\end{equation}

Đặt $n=|V|$ và $p=|\mathcal{R}|$. Chúng tôi định nghĩa các ma trận incidence
cho thân quy tắc
$\mathbf{B}^{+},\mathbf{B}^{-}\in\{0,1\}^{p\times n}$ và các ma trận incidence
cho đầu quy tắc
$\mathbf{H}^{+},\mathbf{H}^{-}\in\{0,1\}^{n\times p}$. Với
$\mathbf{k}=(\mathbf{B}^{+}+\mathbf{B}^{-})\mathbf{1}$, các thân quy tắc và mức
kích hoạt được tính bởi
\begin{equation}
\mathbf{b}
=
\mathbb{I}
[\mathbf{B}^{+}\mathbf{t}
+\mathbf{B}^{-}\mathbf{f}
=\mathbf{k}],
\qquad
\mathbf{a}
=
\mathbf{b}\odot(\mathbf{1}-\mathbf{e}),
\label{eq:rule_activation}
\end{equation}
trong đó $\mathbf{e}$ biểu diễn các ngoại lệ đang hoạt động. Điều kiện bằng
$\mathbf{k}$ yêu cầu mọi literal trong thân quy tắc đều được thỏa mãn, qua đó
biểu diễn phép hội trong từng quy tắc.

Hỗ trợ cho các đầu quy tắc được tổng hợp như sau:
\begin{equation}
\mathbf{d}^{+}
=
\mathbb{I}[\mathbf{H}^{+}\mathbf{a}\geq1],
\qquad
\mathbf{d}^{-}
=
\mathbb{I}[\mathbf{H}^{-}\mathbf{a}\geq1].
\label{eq:head_support}
\end{equation}
Phép đặt ngưỡng biểu diễn phép tuyển giữa các quy tắc được kích hoạt độc lập có
cùng hệ quả. Do đó, $(A\land B)\rightarrow Y$ được mã hóa thành một quy tắc có
hai phần tử trong thân, còn $(A\lor B)\rightarrow Y$ được chuẩn hóa thành
$A\rightarrow Y$ và $B\rightarrow Y$. Biểu diễn incidence này có liên hệ với
lập trình logic trong không gian ma trận và tensor
\cite{sakama2021tensor}.

Trạng thái factual được xác định bằng phép lặp tới điểm bất động:
\begin{equation}
\mathbf{x}^{F}
=
\operatorname{Fix}
(\Phi_{\mathcal{R}},\mathbf{x}^{(0)}).
\label{eq:factual_fixed_point}
\end{equation}
Chỉ có hỗ trợ dương hoặc chỉ có hỗ trợ âm lần lượt tạo trạng thái true hoặc
false; không có hỗ trợ hoặc có hỗ trợ xung đột tạo trạng thái unknown. Runtime
đánh giá các quy tắc thưa bằng phép suy luận ký hiệu thay vì thực hiện phép nhân
ma trận dày. Mặc dù engine hỗ trợ thân quy tắc có nhiều literal, dữ liệu hiện
tại chỉ chứa một tiền đề dương trong mỗi quy tắc được import. Vì vậy, khả năng
xử lý phép hội chưa được đánh giá thực nghiệm trong nghiên cứu này.

\subsection{Chẩn đoán tô-pô và can thiệp cấu trúc}

Gọi $\mathcal{M}(P^*)=\{v_1,\ldots,v_{h-1}\}$ là tập sự kiện trung gian của
đường chính và $m\in\mathcal{M}(P^*)$ là một sự kiện trung gian. Phép chẩn đoán
xóa nút tạo
\begin{equation}
G^{-m}=G[V\setminus\{m\}]
\end{equation}
và kiểm tra
\begin{equation}
\rho_{\mathrm{top}}(s,Y\mid m)
=
\mathbb{I}
[Y\text{ có thể đạt được từ }s\text{ trong }G^{-m}].
\label{eq:topological_ablation}
\end{equation}
Phép toán này chỉ đo tính đạt được tô-pô có giới hạn và không được sử dụng làm
định nghĩa phản thực của LegalSCM.

Với câu hỏi yêu cầu loại bỏ sự kiện trung gian một cách tường minh, LegalSCM
thay vào đó áp dụng can thiệp cấu trúc \cite{pearl2009causality}
\begin{equation}
do(X_m=0),\qquad 0\equiv\mathrm{FALSE}.
\end{equation}
Can thiệp cố định $X_m$ ở trạng thái false và vô hiệu hóa mọi quy tắc xác định
biến này:
\begin{equation}
\mathcal{R}^{do(X_m=0)}
=
\{r\in\mathcal{R}\mid
\operatorname{var}(\operatorname{head}(r))\neq m\}.
\end{equation}
Sau đó, trạng thái phản thực được suy luận lại:
\begin{equation}
\mathbf{x}^{CF(m)}
=
\operatorname{Fix}\!\left(
\Phi_{\mathcal{R}^{do(X_m=0)}},
\mathbf{x}^{(0)}[X_m\leftarrow0]
\right).
\label{eq:cf_fixed_point}
\end{equation}

Sự kiện trung gian $m$ là cần thiết đối với đích $Y$ khi $x_Y^{F}=1$ và
$x_Y^{CF(m)}=0$; không cần thiết khi cả hai giá trị đều bằng $1$; và chưa xác
định trong các trường hợp còn lại. Đây là can thiệp cấu trúc tương đối với mô
hình đã mã hóa, không có bước abduction đối với biến ngoại sinh, và không được
đồng nhất với thao tác xóa $m$ khỏi $G$.

Cuối cùng, hệ thống dự đoán
\begin{equation}
z\in\{
\textsc{supported},
\textsc{reject-direct-claim},
\textsc{uncertain}
\}
\end{equation}
và sinh đầu ra
\begin{equation}
f:(q,G,\mathcal{R})
\rightarrow(P^*,z,\mathcal{E}^*,a,C),
\end{equation}
trong đó $\mathcal{E}^*$ là bằng chứng đã được xác minh, $a$ là câu trả lời và
$C$ là tập căn cứ pháp lý được trích dẫn.

% =========================================================
% IV. LEGALSCM-RAG
% =========================================================
\section{LegalSCM-RAG}
\label{sec:legalscm_rag}

LegalSCM-RAG tích hợp truy xuất đường dẫn nhân quả, xác minh LegalSCM phụ thuộc
câu hỏi và sinh câu trả lời dựa trên bằng chứng. Với câu hỏi $q$, pipeline là
\begin{equation}
\begin{aligned}
q
&\rightarrow
\text{Truy xuất đường dẫn nhân quả}
\rightarrow P^*\\
&\rightarrow
\text{Xác minh phụ thuộc câu hỏi}
\rightarrow (z,\mathcal{E}^*)\\
&\rightarrow
\text{Sinh câu trả lời}
\rightarrow (a,C),
\end{aligned}
\label{eq:overall_pipeline}
\end{equation}
trong đó $P^*$ là đường suy luận được chọn, $z$ là quyết định tương ứng với câu
hỏi, $\mathcal{E}^*$ là bằng chứng cuối cùng, $a$ là câu trả lời và $C$ là tập
căn cứ pháp lý được trích dẫn.

\subsection{Truy xuất đường dẫn nhân quả}

Hệ thống mã hóa câu hỏi thành embedding và truy xuất các sự kiện cùng quy tắc
liên quan từ bộ nhớ nhân quả. Các sự kiện được truy xuất đóng vai trò là các
điểm khởi đầu để mở rộng đồ thị điều kiện--hệ quả có hướng và xây dựng tập
đường ứng viên $\mathcal{P}(q)$. Mỗi đường được xếp hạng bằng các tín hiệu ngữ
nghĩa và cấu trúc:
\begin{equation}
S(P_i,q)=
\alpha S_{\mathrm{sem}}
+\beta S_{\mathrm{graph}}
+\gamma S_{\mathrm{seed}}
+\delta S_{\mathrm{path}}.
\end{equation}
Đường chính được chọn như sau:
\begin{equation}
P^*=\arg\max_{P_i\in\mathcal{P}(q)}S(P_i,q).
\end{equation}
Khác với RAG thông thường, ngữ cảnh được truy xuất bao gồm các sự kiện pháp lý,
quy tắc và quan hệ có hướng kết nối chúng, thay vì chỉ gồm các đoạn văn độc lập.

\subsection{Xác minh LegalSCM phụ thuộc câu hỏi}

LegalSCM hiện thực hóa dạng ký hiệu thưa tương ứng với biểu diễn incidence
thân/đầu quy tắc ba giá trị ở Mục~\ref{sec:problem_formulation}. Hệ thống lặp
việc áp dụng các quy tắc được truy xuất cho đến khi đạt điểm bất động trên
$\{\mathrm{TRUE},\mathrm{FALSE},\mathrm{UNKNOWN}\}$.

Bộ xác minh chọn một trong ba route:

\begin{enumerate}
    \item một khẳng định trực tiếp $A\rightarrow C$ được đánh giá bằng cách
    kiểm tra xem có cạnh trực tiếp đã được ground giữa hai đầu mút hay không;
    \item một câu hỏi pháp lý thông thường được xác thực factual bằng LegalSCM
    trên đường đã chọn mà không thực hiện can thiệp; và
    \item một câu hỏi nêu rõ việc loại bỏ sự kiện trung gian hoặc hỏi về tính
    cần thiết sẽ kích hoạt $do(X_m=0)$ và so sánh các điểm bất động factual với
    điểm bất động sau can thiệp.
\end{enumerate}

Đối với route thứ ba, $do(X_m=0)$ cố định sự kiện trung gian $m$ ở trạng thái
sai, vô hiệu hóa các cơ chế quy tắc xác định sự kiện này và suy luận lại điểm
bất động của hệ quy tắc. Phép toán đó khác với kiểm tra tính đạt được trong
$G^{-m}=G[V\setminus\{m\}]$. Phép xóa nút chỉ được giữ lại như một chẩn đoán
tô-pô có giới hạn và không bao giờ thay thế can thiệp cấu trúc
\cite{pearl2009causality}. Trạng thái cần thiết thu được chỉ tương đối với mô
hình quy tắc quy phạm đã mã hóa, không phải là ước lượng hiệu ứng nhân quả thống
kê hay quan hệ nhân quả thực tế ở cấp cá thể.

\subsection{Cổng cam kết chọn lọc}

Kết quả của bộ xác minh chỉ trở thành kết quả có thẩm quyền khi các phép kiểm
tra liên quan về grounding, độ phủ đường chính, độ tin cậy và tính nhất quán
trạng thái đều đạt yêu cầu. Chúng tôi định nghĩa
\begin{equation}
G_{\mathrm{commit}}
=
\mathbb{I}
[
g_{\mathrm{ground}}
\land g_{\mathrm{path}}
\land (s_{\mathrm{SCM}}\geq\tau)
\land g_{\mathrm{consistent}}
].
\end{equation}
Quyết định và bằng chứng cuối cùng sau đó được chọn bởi
\begin{equation}
z=
\begin{cases}
z_{\mathrm{SCM}}, & G_{\mathrm{commit}}=1,\\
z_{\mathrm{LLM}}, & G_{\mathrm{commit}}=0,
\end{cases}
\qquad
\mathcal{E}^*=
\begin{cases}
\mathcal{E}_{\mathrm{verified}}, & G_{\mathrm{commit}}=1,\\
\mathcal{E}_{\mathrm{path}}, & G_{\mathrm{commit}}=0.
\end{cases}
\label{eq:selective_commit}
\end{equation}

Khi cổng không đạt, hệ thống không chuyển một chẩn đoán tô-pô thành quyết định
SCM và cũng không tự động ép kết quả thành \textsc{uncertain}. Thay vào đó, hệ
thống giữ nguyên bằng chứng ban đầu của Causal Path RAG và để mô hình sinh dùng
chung đưa ra quyết định. Do đó, Full LegalSCM là một hệ thống lai chọn lọc chứ
không phải một bộ phân loại mà SCM có thẩm quyền trên mọi câu hỏi.

\subsection{Sinh câu trả lời và khả năng kiểm toán}

Mô hình ngôn ngữ nhận đầu vào
\begin{equation}
I_q=(q,P^*,z,\mathcal{E}^*)
\end{equation}
và sinh
\begin{equation}
(a,C)=f_{\theta}(I_q).
\end{equation}
Tất cả các hệ thống được so sánh đều sử dụng cùng schema sinh và cùng tham số
giải mã. LegalSCM-RAG còn ghi lại route được chọn, kết quả grounding, trạng
thái can thiệp, quyết định commit, trạng thái fallback và bằng chứng cuối cùng.
Sự phân tách này cho phép kiểm toán độc lập quá trình truy xuất, suy luận quy
tắc, can thiệp cấu trúc và sinh câu trả lời.

% =========================================================
% V. EXPERIMENTAL SETUP
% =========================================================
\section{Thiết lập thực nghiệm}
\label{sec:experimental_setup}

\subsection{Dữ liệu và các hệ thống so sánh}

Chúng tôi đánh giá trên 250 câu hỏi nhiều bước được xây dựng từ Bộ luật Hình sự
Việt Nam \cite{blhs2025} và đồ thị pháp luật dẫn xuất. Trong số đó, 236 câu có
đường chuẩn tuyến tính. Các nhãn quyết định gồm 230 mẫu
\textsc{supported} và 20 mẫu \textsc{reject-direct-claim}, không có mẫu chuẩn
\textsc{uncertain}. Bộ định tuyến xử lý 13 mẫu âm bằng phép kiểm tra cạnh trực
tiếp và 7 mẫu loại bỏ sự kiện trung gian tường minh bằng $do(X_m=0)$. Các
trường chuẩn và trường đánh giá không bao giờ được truy cập trong quá trình sinh
dự đoán. Do các câu hỏi và đồ thị truy xuất có cùng nguồn dữ liệu có cấu trúc,
thực nghiệm chủ yếu đo tính hợp lệ nội tại trong môi trường có kiểm soát.

Chúng tôi so sánh: (1) LLM closed-book cho phép abstention; (2) LLM closed-book
buộc quyết định nhị phân; (3) Vanilla RAG; (4) Causal Path RAG; và (5) Full
LegalSCM. Bốn hệ thống cho phép abstention sử dụng cùng bộ sinh, quy tắc prompt,
schema đầu ra, ngân sách ngữ cảnh và tham số giải mã. Biến thể buộc quyết định
nhị phân chỉ thay đổi chính sách quyết định và được xem là một đối chứng phân
tích độ nhạy. Không hệ thống nào sử dụng fallback trích xuất cho câu trả lời.

\subsection{Cấu hình}

Qwen3-8B được chạy qua Ollama, trong khi BGE-M3
\cite{chen2024bgem3} và FAISS đảm nhiệm truy xuất ngữ nghĩa.

\begin{table}[t]
\centering
\caption{Cấu hình thực nghiệm chính.}
\label{tab:experimental_config}
\small
\begin{tabular}{lc}
\hline
\textbf{Tham số} & \textbf{Giá trị}\\
\hline
Bộ sinh / mô hình embedding & Qwen3-8B / BGE-M3\\
Nhiệt độ / seed & 0.1 / 42\\
Token đầu ra / cửa sổ ngữ cảnh & 1200 / 32768\\
Số ký tự ngữ cảnh tối đa & 18000\\
Top-$K$ Vanilla / sự kiện / quy tắc & 5 / 8 / 8\\
Top-$K$ bằng chứng nhân quả cuối cùng & 12\\
Số hop đồ thị tối đa & 2\\
Ngưỡng commit / grounding & 0.55 / 0.80\\
\hline
\end{tabular}
\end{table}

\subsection{Chỉ số và kiểm định thống kê}

Chúng tôi báo cáo Decision Accuracy, Balanced Accuracy và Macro-F1. Do không có
mẫu chuẩn nào thuộc lớp \textsc{uncertain}, hai chỉ số sau được lấy trung bình
trên hai lớp có trong dữ liệu chuẩn. Token-F1 và ROUGE-L đánh giá câu trả lời,
trong khi F1 dựa trên tập hợp đánh giá căn cứ được trích dẫn. Các thành phần
truy xuất phù hợp được đánh giá bằng Recall@5, MRR và MAP; Top-1 Exact Path và
Oracle Exact Path lần lượt đánh giá đường được chọn và đường ứng viên. Các điểm
Oracle chỉ dùng cho mục đích chẩn đoán. Dự đoán bị thiếu hoặc thất bại nhận điểm
0 trên các chỉ số end-to-end dùng chung, còn thời gian chạy được lấy trung bình
trên các dự đoán thành công.

Mọi so sánh đều theo cặp trên cùng 250 câu hỏi. Chúng tôi dùng 5.000 mẫu
bootstrap cho khoảng tin cậy 95\%, kiểm định McNemar chính xác cho Decision
Accuracy và 5.000 phép đổi dấu ngẫu nhiên theo cặp cho Answer F1, ROUGE-L và
Citation F1, cùng hiệu chỉnh Holm. Chúng tôi cũng kiểm toán route xác minh,
quyết định commit, fallback, và việc $do(X_m=0)$ có được yêu cầu và thực thi hay
không.

% =========================================================
% VI. RESULTS AND ANALYSIS
% =========================================================
\section{Kết quả và phân tích}
\label{sec:results_analysis}

\subsection{Kết quả tổng thể và các đối chứng closed-book}

Cả năm lượt chạy đều sinh đủ 250/250 dự đoán mà không có lỗi pipeline.
Bảng~\ref{tab:main_results} trình bày kết quả end-to-end.

\begin{table}[t]
\centering
\caption{Hiệu năng trên 250 câu hỏi (đơn vị phần trăm, trừ thời gian).}
\label{tab:main_results}
\scriptsize
\setlength{\tabcolsep}{1.35pt}
\begin{tabular}{@{}lrrrrrrr@{}}
\toprule
\textbf{Phương pháp}
& \textbf{Acc.}
& \textbf{B.Acc.}
& \textbf{M-F1}
& \textbf{Ans.F1}
& \textbf{R-L}
& \textbf{Cit.F1}
& \textbf{Giây}\\
\midrule
Closed-book (abstain) & 4.40 & 27.50 & 22.00 & 15.08 & 13.19 & 0.00 & 1.031\\
Closed-book (forced) & 9.20 & 50.65 & 8.78 & 18.44 & 16.19 & 0.27 & 0.900\\
Vanilla RAG & 55.20 & 71.09 & 49.88 & 28.18 & 25.68 & 50.75 & 2.369\\
Causal Path RAG & 88.80 & 89.35 & 76.30 & 37.27 & 36.40 & 58.41 & 1.947\\
Full LegalSCM & \textbf{96.80} & \textbf{98.26} & \textbf{91.67}
& \textbf{37.60} & \textbf{36.92} & \textbf{58.89} & 3.238\\
\bottomrule
\end{tabular}
\end{table}

Điểm closed-book thấp không bắt nguồn từ lượt chạy thất bại hay lỗi phân tích
đầu ra. Khi cho phép abstention, mô hình trả \textsc{uncertain} cho 220 câu
hỏi; do tập chuẩn không chứa nhãn \textsc{uncertain}, chỉ 11 dự đoán là đúng.
Khi bị buộc đưa ra quyết định nhị phân, mô hình dự đoán
\textsc{reject-direct-claim} cho 247 trong số 250 câu hỏi và đạt độ chính xác
9,20\%. Như vậy, loại bỏ abstention tạo ra thiên lệch bác bỏ gần như đồng nhất
thay vì khôi phục năng lực pháp luật. Vì một bộ phân loại luôn dự đoán
\textsc{supported} sẽ đạt độ chính xác thô 92\% trên benchmark lệch lớp này,
Balanced Accuracy và Macro-F1 cung cấp nhiều thông tin hơn so với chỉ riêng
Accuracy.

\subsection{Ảnh hưởng của xác minh chọn lọc}

Full LegalSCM trả lời đúng 242 câu hỏi, so với 222 câu của Causal Path RAG. Chỉ
Full đúng trên 23 mẫu, chỉ Causal Path RAG đúng trên 3 mẫu, cả hai cùng đúng
trên 219 mẫu và cả hai cùng sai trên 5 mẫu. Mức cải thiện 8,0 điểm thu được có
khoảng tin cậy bootstrap theo cặp 95\% là $[4{,}4;12{,}0]$, với kiểm định
McNemar chính xác $p=8{,}8\times10^{-5}$ và $p$ hiệu chỉnh Holm $=0{,}0011$.
Macro-F1 tăng từ 76,30\% lên 91,67\%, trong khi F1 của lớp âm tăng từ 59,02\%
lên 85,11\%.

Các khác biệt về sinh văn bản tự do nhỏ hơn đáng kể: Answer F1, ROUGE-L và
Citation F1 lần lượt thay đổi $+0{,}33$, $+0{,}52$ và $+0{,}47$ điểm. Khoảng
tin cậy của cả ba chỉ số đều chứa 0 và cả ba giá trị $p$ hiệu chỉnh Holm đều
bằng 1,0. Do đó, mức cải thiện đo được chỉ đặc trưng cho độ chính xác của quyết
định có cấu trúc.

\subsection{Kiểm toán xác minh và lỗi}

Các thống kê về route và commit của Full LegalSCM được trình bày trong
Bảng~\ref{tab:route_audit}.

\begin{table}[t]
\centering
\caption{Kiểm toán xác minh của Full LegalSCM.}
\label{tab:route_audit}
\small
\begin{tabular}{lr}
\toprule
\textbf{Đại lượng kiểm toán} & \textbf{Số lượng}\\
\midrule
Route xác thực factual & 230\\
Route cạnh trực tiếp & 13\\
Route can thiệp sự kiện trung gian & 7\\
Bộ xác minh đã commit & 230\\
Fallback do không có đường chính & 20\\
Đã thực thi $do(X_m=0)$ & 7\\
\bottomrule
\end{tabular}
\end{table}

Cả 23 trường hợp chỉ Full cải thiện đều xuất hiện sau khi bộ xác minh commit:
21 trường hợp đến từ xác thực factual và 2 trường hợp đến từ kiểm tra cạnh trực
tiếp. Không trường hợp nào đến từ route can thiệp. Trong 20 trường hợp
fallback, Full không có cải thiện riêng, còn chỉ Causal Path RAG đúng trên 3
mẫu.

Cả tám lỗi của Full đều xảy ra sau fallback do không có đường chính; bảy trong
số đó là câu hỏi suy luận ngược. Không có lỗi nào xảy ra sau một kết quả đã
được bộ xác minh commit. Trong 20 mẫu âm, Full trả lời đúng cả 13 trường hợp
cạnh trực tiếp và cả 7 trường hợp can thiệp sự kiện trung gian, so với độ chính
xác tổng thể 90\% của Causal Path RAG. Tuy nhiên, Causal Path RAG cũng trả lời
đúng cả 7 trường hợp can thiệp. Vì vậy, các can thiệp đã thực thi có thể được
kiểm toán nhưng không tạo lợi thế theo cặp riêng biệt trên tập con này.

\subsection{Truy xuất, xếp hạng đường và thời gian chạy}

Causal Path RAG và Full sử dụng đầu ra truy xuất giống hệt nhau. Article,
Rule và Event Recall@5 của chúng lần lượt là 75,40\%, 71,20\% và 74,79\%, trong
khi Vanilla RAG đạt Article Recall@5 là 60,07\%. Top-1 Exact Path đạt 38,56\%,
so với Oracle Exact Path 75,42\%; khoảng cách này cho thấy một đường phù hợp
thường đã được truy xuất nhưng chưa được xếp hạng đầu tiên.

Full cần 3,238 giây cho mỗi câu hỏi, so với 1,947 giây của Causal Path RAG. Vì
vậy, xác minh chọn lọc cải thiện 8,0 điểm độ chính xác quyết định với thời gian
chạy tăng thêm xấp xỉ 66\%.

% =========================================================
% VII. DISCUSSION AND LIMITATIONS
% =========================================================
\section{Thảo luận và hạn chế}
\label{sec:discussion_limitations}

Full LegalSCM đạt độ chính xác quyết định 96,80\%, so với 88,80\% của Causal
Path RAG. Do hai hệ thống sử dụng đầu ra truy xuất giống hệt nhau và cả 23
trường hợp chỉ Full cải thiện đều xuất hiện sau khi bộ xác minh commit, mức cải
thiện được khu trú ở bước xác thực factual đã commit và bước kiểm tra cạnh trực
tiếp có grounding. Tuy nhiên, Answer F1, ROUGE-L và Citation F1 không cải thiện
có ý nghĩa, cho thấy tính đúng đắn của quyết định có cấu trúc và chất lượng câu
trả lời tự do vẫn là hai mục tiêu riêng biệt.

Kết quả closed-book cần được diễn giải thận trọng. Độ chính xác thấp của chúng
phản ánh hợp đồng nhãn định hướng xác minh, việc không có bằng chứng pháp luật
được truy xuất và sự mất cân bằng lớp lớn, thay vì năng lực pháp luật tổng quát
của Qwen3-8B. Điều này tạo động lực để báo cáo Balanced Accuracy và Macro-F1
cùng với Accuracy thô.

Thực nghiệm thực thi $do(X_m=0)$ trên bảy câu hỏi loại bỏ sự kiện trung gian,
nhưng Causal Path RAG cũng trả lời đúng cả bảy câu. Vì vậy, các kết quả hiện tại
chưa tách biệt được lợi ích của can thiệp cấu trúc. Tương tự, mặc dù engine hỗ
trợ thân quy tắc có quan hệ hội và các cơ chế tuyển thay thế, các quy tắc được
nhập chỉ chứa một tiền đề dương trên mỗi bản ghi. Cả suy luận can thiệp và suy
luận nhiều tiền đề đều cần các benchmark chuyên biệt lớn hơn.

Truy xuất cũng vẫn là một nút thắt. Top-1 Exact Path đạt 38,56\%, so với Oracle
Exact Path 75,42\%, và cả tám lỗi của Full đều xảy ra sau fallback do thiếu
đường chính. Hơn nữa, benchmark có quy mô nhỏ, được sinh từ đồ thị và được xây
dựng từ chính nguồn dữ liệu có cấu trúc mà hệ thống sử dụng. Cần có các câu hỏi
pháp lý được chuyên gia rà soát và nằm ngoài đồ thị để đánh giá tính hợp lệ bên
ngoài.

Cuối cùng, LegalSCM là một mô hình quy tắc quy phạm tất định, không phải là một
phương pháp khám phá nhân quả hay ước lượng hiệu ứng thống kê. Ngữ nghĩa can
thiệp của mô hình chỉ tương đối với các cơ chế pháp lý đã mã hóa và không bao
gồm bước abduction biến ngoại sinh cần thiết cho phân tích quan hệ nhân quả
thực tế ở cấp cá thể.

% =========================================================
% VIII. CONCLUSION
% =========================================================
\section{Kết luận}
\label{sec:conclusion}

LegalSCM-RAG kết hợp truy xuất đường dẫn nhân quả, suy luận quy tắc ba giá trị
và xác minh cấu trúc chọn lọc cho hỏi đáp pháp luật tiếng Việt đa bước. Trên 250
câu hỏi, hệ thống cải thiện độ chính xác quyết định 8,0 điểm so với Causal Path
RAG với bằng chứng thống kê theo cặp, trong khi mức cải thiện câu trả lời và căn
cứ trích dẫn không có ý nghĩa.

Các kết quả ủng hộ cơ chế xác minh factual và xác minh khẳng định trực tiếp có
chọn lọc trong môi trường có kiểm soát này, nhưng chưa xác lập lợi ích độc lập
của $do(X_m=0)$. Nghiên cứu tương lai sẽ mở rộng các benchmark về can thiệp sự
kiện trung gian và quy tắc nhiều tiền đề, cải thiện xếp hạng đường cùng
Grounding, áp đặt tính nhất quán giữa câu trả lời và quyết định, đồng thời tổ
chức đánh giá chuyên gia trên các câu hỏi pháp lý không được sinh từ đồ thị
nguồn.

% =========================================================
% REFERENCES
% =========================================================
\begin{thebibliography}{99}

\bibitem{lewis2020retrieval}
P. Lewis, E. Perez, A. Piktus, et al.,
``Retrieval-Augmented Generation for Knowledge-Intensive NLP Tasks,''
in \textit{Advances in Neural Information Processing Systems},
vol. 33, 2020.

\bibitem{karpukhin2020dense}
V. Karpukhin, B. Oguz, S. Min, et al.,
``Dense Passage Retrieval for Open-Domain Question Answering,''
in \textit{Proceedings of EMNLP}, 2020.

\bibitem{yang2018hotpotqa}
Z. Yang, P. Qi, S. Zhang, Y. Bengio, W. W. Cohen,
R. Salakhutdinov, and C. D. Manning,
``HotpotQA: A Dataset for Diverse, Explainable Multi-hop
Question Answering,''
in \textit{Proceedings of EMNLP}, 2018.

\bibitem{trivedi2022musique}
H. Trivedi, N. Balasubramanian, T. Khot, and A. Sabharwal,
``MuSiQue: Multihop Questions via Single-hop Question Composition,''
\textit{Transactions of the Association for Computational Linguistics},
vol. 10, 2022.

\bibitem{edge2024graphrag}
D. Edge, H. Trinh, N. Cheng, et al.,
``From Local to Global: A Graph RAG Approach to
Query-Focused Summarization,''
\textit{arXiv preprint arXiv:2404.16130}, 2024.

\bibitem{wang2025causalrag}
N. Wang, X. Han, J. Singh, J. Ma, and V. Chaudhary,
``CausalRAG: Integrating Causal Graphs into
Retrieval-Augmented Generation,''
in \textit{Findings of the Association for Computational Linguistics:
ACL 2025}, pp. 22680--22693, 2025.

\bibitem{guo2023counterfactual}
W. Guo, Q. Gong, Y. Rao, and H. Lai,
``Counterfactual Multihop QA: A Cause-Effect Approach for
Reducing Disconnected Reasoning,''
in \textit{Proceedings of the 61st Annual Meeting of the
Association for Computational Linguistics},
pp. 4214--4226, 2023.

\bibitem{chalkidis2022lexglue}
I. Chalkidis, A. Jana, D. Hartung, et al.,
``LexGLUE: A Benchmark Dataset for Legal Language
Understanding in English,''
in \textit{Proceedings of the 60th Annual Meeting of the
Association for Computational Linguistics},
pp. 4310--4330, 2022.

\bibitem{pearl2009causality}
J. Pearl,
\textit{Causality: Models, Reasoning, and Inference},
2nd ed. Cambridge University Press, 2009.

\bibitem{chen2024bgem3}
J. Chen, S. Xiao, P. Zhang, K. Luo, D. Lian, and Z. Liu,
``BGE M3-Embedding: Multi-Lingual, Multi-Functionality,
Multi-Granularity Text Embeddings Through Self-Knowledge Distillation,''
\textit{arXiv preprint arXiv:2402.03216}, 2024.

\bibitem{blhs2025}
Văn phòng Quốc hội,
\textit{Văn bản hợp nhất số 135/VBHN-VPQH -- Bộ luật Hình sự},
Hà Nội, ngày 05 tháng 9 năm 2025.

% ---------------------------------------------------------
% Tài liệu tham khảo bổ sung [12]--[20]
% ---------------------------------------------------------

\bibitem{gao2024ragsurvey}
Y. Gao, Y. Xiong, X. Gao, et al.,
``Retrieval-Augmented Generation for Large Language Models: A Survey,''
\textit{arXiv preprint arXiv:2312.10997}, 2024.

\bibitem{izacard2021fid}
G. Izacard and E. Grave,
``Leveraging Passage Retrieval with Generative Models for Open Domain
Question Answering,''
in \textit{Proceedings of the 16th Conference of the European Chapter of the
Association for Computational Linguistics (EACL)},
pp. 874--880, 2021.

\bibitem{asai2024selfrag}
A. Asai, Z. Wu, Y. Wang, A. Sil, and H. Hajishirzi,
``Self-RAG: Learning to Retrieve, Generate, and Critique through
Self-Reflection,''
in \textit{International Conference on Learning Representations (ICLR)}, 2024.

\bibitem{sarthi2024raptor}
P. Sarthi, S. Abdullah, A. Tuli, S. Khanna, A. Goldie, and C. D. Manning,
``RAPTOR: Recursive Abstractive Processing for Tree-Organized Retrieval,''
in \textit{International Conference on Learning Representations (ICLR)}, 2024.

\bibitem{sun2024thinkongraph}
J. Sun, C. Xu, L. Tang, et al.,
``Think-on-Graph: Deep and Responsible Reasoning of Large Language Model on
Knowledge Graph,''
in \textit{International Conference on Learning Representations (ICLR)}, 2024.

\bibitem{parekh2025ccrag}
J. R. Parekh, P. Jiang, and J. Han,
``CC-RAG: Structured Multi-Hop Reasoning via Theme-Based Causal Graphs,''
\textit{arXiv preprint arXiv:2506.08364}, 2025.

\bibitem{guha2023legalbench}
N. Guha, J. Nyarko, D. E. Ho, et al.,
``LegalBench: A Collaboratively Built Benchmark for Measuring Legal Reasoning
in Large Language Models,''
in \textit{Advances in Neural Information Processing Systems},
vol. 36, 2023.

\bibitem{zheng2021casehold}
L. Zheng, N. Guha, B. R. Anderson, P. Henderson, and D. E. Ho,
``When Does Pretraining Help? Assessing Self-Supervised Learning for Law and
the CaseHOLD Dataset of 53,000+ Legal Holdings,''
in \textit{Proceedings of the 18th International Conference on Artificial
Intelligence and Law (ICAIL)},
pp. 159--168, 2021.

\bibitem{chalkidis2020legalbert}
I. Chalkidis, M. Fergadiotis, P. Malakasiotis, N. Aletras, and
I. Androutsopoulos,
``LEGAL-BERT: The Muppets straight out of Law School,''
in \textit{Findings of the Association for Computational Linguistics:
EMNLP 2020},
pp. 2898--2904, 2020.

\bibitem{kleene1952metamathematics}
S. C. Kleene,
\textit{Introduction to Metamathematics}.
North-Holland, Amsterdam, 1952.

\bibitem{fitting1985kripke}
M. Fitting,
``A Kripke--Kleene Semantics for Logic Programs,''
\textit{The Journal of Logic Programming},
vol. 2, no. 4, pp. 295--312, 1985.

\bibitem{sato2020threevalued}
T. Sato, C. Sakama, and K. Inoue,
``From 3-valued Semantics to Supported Model Computation for
Logic Programs in Vector Spaces,''
in \textit{Proceedings of the 12th International Conference on
Agents and Artificial Intelligence (ICAART), Volume 2},
pp. 758--765, SciTePress, 2020.
doi: 10.5220/0009093407580765.

\bibitem{sakama2021tensor}
C. Sakama, K. Inoue, and T. Sato,
``Logic Programming in Tensor Spaces,''
\textit{Annals of Mathematics and Artificial Intelligence},
vol. 89, no. 12, pp. 1133--1153, 2021.


\end{thebibliography}

\bibliographystyle{splncs04}

\end{document}
